Repeated surveys produce estimated distributions for many populations, and those
estimates are increasingly used to calibrate a prediction band for a further
population. The calibration objects then carry sampling error of their own, so a
band calibrated on them covers a further survey estimate but not automatically
the underlying population distribution. We separate these two targets and ask
when the sampling error can be removed. A conformal band treating populations as
exchangeable gives finite-sample coverage for a further survey estimate;
coverage for the underlying distribution requires either a sampling-error bound,
which is shape-free and wide, or a scale correction, which requires the
standardised observed and underlying curves to share a law. We then show that
the two conditions an adaptive implementation of that correction gates on are
not independent. The algebra is classical, from variance-component estimation,
but its consequence for a criterion deciding whether a correction may be used
has not been drawn out. The correction shrinks the target scale in proportion to
the design share and its reliability diagnostic is evaluated on that shrunken
scale, so the requirement is a closed-form boundary rising steeply in the design
share. A reported floor of 94 populations is that boundary
at a design share of zero, where nothing needs removing; where the need gate
opens it is 153. Across 54 configurations of European Social Survey estimates it
predicts the gate correctly in 53. Where it activates it is 15 to 19 percent
narrower than the uncorrected band. Refining the calibration unit can reach those
configurations, but only along a dimension on which the domains genuinely
differ: cutting the survey by region comes within reach of the requirement,
while cutting it by sex falls an order of magnitude short.
